\section{Roadmap Update after Experiment 11B}
\label{sec:roadmap-post-11b}

Experiment~11B passes the practical-certificate gate sufficiently strongly that the project should now leave the quadratic family. The calibrated rule is no longer only an oracle construction: on the strongly rotated benchmark a native-coordinate certificate with sufficiently frequent dense calibration achieves zero measured harmful applied events, substantially improves the tuned no-certificate event optimizer, and includes a point that dominates the selected EF-TopK reference in tail error, whole-run objective, and total payload bits.

At the same time, the experiment identifies two limitations that must be carried forward rather than hidden:
\begin{enumerate}
  \item \textbf{Calibration is the dominant communication cost.} At the strongest native rotated operating point ($C=25$), roughly 97\% of payload bits come from dense calibration gradients rather than the sparse event stream.
  \item \textbf{Certificate tightness depends strongly on representation geometry.} Sparse calibration is almost ideal in the diagonal/basis-aligned problem but must be much more frequent in raw rotated coordinates.
\end{enumerate}
The next experiment should therefore test whether the complete mechanism survives on a nonquadratic learning objective where smoothness bounds can be constructed from real data statistics and where model quality, rather than distance to a known quadratic optimum, is the primary metric.

\subsection{Recommended next step: Experiment 12A -- federated logistic regression}

The next study should be \textbf{Experiment~12A: non-IID federated logistic regression with calibrated sparse descent events}. Logistic regression is the correct intermediate benchmark because it introduces a genuine nonlinear loss while retaining convexity and usable global curvature bounds. It is substantially more informative than another synthetic quadratic but still simple enough that failures can be diagnosed mathematically.

For client $i$, use the regularized binary logistic objective
\begin{equation}
  F_i(w)
  =
  \frac{1}{n_i}\sum_{m=1}^{n_i}
  \log\!\left(1+\exp(-y_{im}x_{im}^\top w)\right)
  +\frac{\lambda}{2}\norm{w}^2.
\end{equation}
The Hessian satisfies
\begin{equation}
  \nabla^2F_i(w)
  =
  \frac{1}{n_i}X_i^\top D_i(w)X_i+\lambda I,
  \qquad
  0\preceq D_i(w)\preceq\frac14 I.
\end{equation}
Therefore clients can construct a state-independent componentwise bound
\begin{equation}
  M_{i,jk}
  =
  \frac{1}{4n_i}\sum_m\abs{x_{im,j}x_{im,k}}
  +\lambda\mathbf 1_{j=k},
\end{equation}
which guarantees
\begin{equation}
  \abs{[\nabla^2F_i(w)]_{jk}}
  \leq M_{i,jk}.
\end{equation}
This is the direct nonquadratic analogue of the Experiment~11B componentwise drift envelope and makes logistic regression an unusually clean next test.

\subsection{The new communication question}

Experiment~12A must count not only event and recurring gradient-calibration bits, but also any curvature-summary information required to build the certificate. A full $d\times d$ matrix $M_i$ is acceptable as an oracle/upper-information reference for small $d$, but it is not a scalable solution. The experiment should therefore compare at least:
\begin{enumerate}
  \item a full componentwise curvature matrix $M_i$;
  \item a diagonal-only curvature bound;
  \item a row-norm or block-diagonal bound;
  \item the no-certificate global $q(t)$ event baseline;
  \item EF-TopK/error-feedback sparsification;
  \item full-precision asynchronous SGD.
\end{enumerate}
The one-time bits required to transmit or configure these curvature summaries must be reported explicitly. This turns the dominant limitation discovered in 11B into a measurable systems-design trade-off rather than postponing it.

\subsection{Data heterogeneity}

The first logistic experiment should remain controlled. Use a public or synthetic binary-classification dataset partitioned across approximately 10--20 clients with tunable label/feature heterogeneity. At minimum include:
\begin{itemize}
  \item an approximately IID split;
  \item moderate Dirichlet label skew;
  \item a stronger non-IID split that produces meaningful local/global gradient disagreement.
\end{itemize}
The primary comparison should not rely on a single favorable partition. Fixed seeds and common stochastic mini-batch draws should be used where possible.

\subsection{Primary metrics}

The main plots should now become learning-system plots rather than quadratic-only diagnostics:
\begin{equation}
  \boxed{\text{test loss / test accuracy vs cumulative transmitted bits}.}
\end{equation}
Also report training objective, calibration count, accepted/candidate event ratio, harmful-event fraction whenever the full objective can be evaluated, client participation, and wall-clock/local-gradient evaluations. Communication should be decomposed into
\begin{equation}
  B_{\rm total}
  =
  B_{\rm events}
  +B_{\rm gradient\;calibration}
  +B_{\rm curvature\;summary}.
\end{equation}
This decomposition is essential because 11B showed that the event stream can look extremely efficient while calibration silently dominates the total budget.

\subsection{Experiment 12A decision gate}

The calibrated-event approach should proceed toward a neural network only if a practical curvature-summary variant produces a useful test-performance--bits Pareto point relative to EF-TopK and full-precision training. A full-matrix oracle certificate alone is not sufficient justification for scaling.

If the full componentwise certificate works but diagonal/block summaries fail, the next step should be a representation/preconditioning study before neural networks. This would connect directly to the Experiment~10D basis result. If even the full componentwise logistic certificate fails, the descent-certificate branch should be reconsidered before further scaling.

If a diagonal or small-block certificate succeeds, the project can proceed to the originally planned small-neural-network experiment, where certificates would be formed per layer/block rather than per raw parameter.

The updated progression is therefore
\begin{equation}
  \boxed{12\mathrm{A}\;\text{logistic regression + practical certificate}\;\rightarrow\;12\mathrm{B}\;\text{representation/block refinement if needed}\;\rightarrow\;13\;\text{small neural network}.}
\end{equation}
